% generated by GAPDoc2LaTeX from XML source (Frank Luebeck)
\documentclass[a4paper,11pt]{report}

\usepackage[top=37mm,bottom=37mm,left=27mm,right=27mm]{geometry}
\sloppy
\pagestyle{myheadings}
\usepackage{amssymb}
\usepackage[utf8]{inputenc}
\usepackage{makeidx}
\makeindex
\usepackage{color}
\definecolor{FireBrick}{rgb}{0.5812,0.0074,0.0083}
\definecolor{RoyalBlue}{rgb}{0.0236,0.0894,0.6179}
\definecolor{RoyalGreen}{rgb}{0.0236,0.6179,0.0894}
\definecolor{RoyalRed}{rgb}{0.6179,0.0236,0.0894}
\definecolor{LightBlue}{rgb}{0.8544,0.9511,1.0000}
\definecolor{Black}{rgb}{0.0,0.0,0.0}

\definecolor{linkColor}{rgb}{0.0,0.0,0.554}
\definecolor{citeColor}{rgb}{0.0,0.0,0.554}
\definecolor{fileColor}{rgb}{0.0,0.0,0.554}
\definecolor{urlColor}{rgb}{0.0,0.0,0.554}
\definecolor{promptColor}{rgb}{0.0,0.0,0.589}
\definecolor{brkpromptColor}{rgb}{0.589,0.0,0.0}
\definecolor{gapinputColor}{rgb}{0.589,0.0,0.0}
\definecolor{gapoutputColor}{rgb}{0.0,0.0,0.0}

%%  for a long time these were red and blue by default,
%%  now black, but keep variables to overwrite
\definecolor{FuncColor}{rgb}{0.0,0.0,0.0}
%% strange name because of pdflatex bug:
\definecolor{Chapter }{rgb}{0.0,0.0,0.0}
\definecolor{DarkOlive}{rgb}{0.1047,0.2412,0.0064}


\usepackage{fancyvrb}

\usepackage{mathptmx,helvet}
\usepackage[T1]{fontenc}
\usepackage{textcomp}


\usepackage[
            pdftex=true,
            bookmarks=true,        
            a4paper=true,
            pdftitle={Written with GAPDoc},
            pdfcreator={LaTeX with hyperref package / GAPDoc},
            colorlinks=true,
            backref=page,
            breaklinks=true,
            linkcolor=linkColor,
            citecolor=citeColor,
            filecolor=fileColor,
            urlcolor=urlColor,
            pdfpagemode={UseNone}, 
           ]{hyperref}

\newcommand{\maintitlesize}{\fontsize{50}{55}\selectfont}

% write page numbers to a .pnr log file for online help
\newwrite\pagenrlog
\immediate\openout\pagenrlog =\jobname.pnr
\immediate\write\pagenrlog{PAGENRS := [}
\newcommand{\logpage}[1]{\protect\write\pagenrlog{#1, \thepage,}}
%% were never documented, give conflicts with some additional packages

\newcommand{\GAP}{\textsf{GAP}}

%% nicer description environments, allows long labels
\usepackage{enumitem}
\setdescription{style=nextline}

%% depth of toc
\setcounter{tocdepth}{1}





%% command for ColorPrompt style examples
\newcommand{\gapprompt}[1]{\color{promptColor}{\bfseries #1}}
\newcommand{\gapbrkprompt}[1]{\color{brkpromptColor}{\bfseries #1}}
\newcommand{\gapinput}[1]{\color{gapinputColor}{#1}}


\begin{document}

\logpage{[ 0, 0, 0 ]}
\begin{titlepage}
\mbox{}\vfill

\begin{center}{\maintitlesize \textbf{ fgab \mbox{}}}\\
\vfill

\hypersetup{pdftitle= fgab }
\markright{\scriptsize \mbox{}\hfill  fgab  \hfill\mbox{}}
{\Huge \textbf{ for finitely generated abelian groups \mbox{}}}\\
\vfill

{\Huge  0.1 \mbox{}}\\[1cm]
{ 25 February 2026 \mbox{}}\\[1cm]
\mbox{}\\[2cm]
{\Large \textbf{ Tian Yuan  \mbox{}}}\\
\hypersetup{pdfauthor= Tian Yuan  }
\end{center}\vfill

\mbox{}\\
{\mbox{}\\
\small \noindent \textbf{ Tian Yuan  }  Email: \href{mailto://ttiiddee@foxmail.com} {\texttt{ttiiddee@foxmail.com}}}\\
\end{titlepage}

\newpage\setcounter{page}{2}
\newpage

\def\contentsname{Contents\logpage{[ 0, 0, 1 ]}}

\tableofcontents
\newpage

     
\chapter{\textcolor{Chapter }{Introduction}}\label{Chapter_Introduction}
\logpage{[ 1, 0, 0 ]}
\hyperdef{L}{X7DFB63A97E67C0A1}{}
{
  

 This package is developed for computation problems of \texttt{FGAb}, i.e. the abbreviation of finitely generated abelian groups. For example, we
are able to compute the kernel and cokernel of a \texttt{FGAbHomomorphism}, along with the direct sum, we can compute all finite limits and colimits, in
particular, pullback and pushout. 

 The basic idea is to use a matrix $f$ over integers to represent a \texttt{FGAb}. The finitely generated abelian group $A$ it represents is equivalent to the cokernel of $f:\mathbb{Z}^n\to \mathbb{Z}^m$. The readers who are familiar with the homological algebra could think $f$ as the free resolutoin of $A$. Actually, almost all algorithms in this package are developed using the
knowledge of homological algebra. For example, we use the mapping cone and
cocone to compute kernel and cokernel. For more details, see the appendix of
our paper $\color{red}{\textrm{(add this later)}}$. 

 }

   
\chapter{\textcolor{Chapter }{\texttt{FGAb}}}\label{Chapter_`FGAb`}
\logpage{[ 2, 0, 0 ]}
\hyperdef{L}{X802A67727D87EFF5}{}
{
  

 The matrices and vectors used in our data structures are \texttt{MatrixObj} and \texttt{VectorObj}. The reason for using them is that we often encounter
0\texttt{\symbol{45}}dimensional matrices and vectors. However, you don't need
to be worried about it as a user, since most functions support naive matrices
and vectors. 

 
\section{\textcolor{Chapter }{\texttt{FGAbElement}}}\label{Chapter_`FGAb`_Section_`FGAbElement`}
\logpage{[ 2, 1, 0 ]}
\hyperdef{L}{X833081227C05A75C}{}
{
  

 

\subsection{\textcolor{Chapter }{FGAbElement (for IsFGAb, IsVector)}}
\logpage{[ 2, 1, 1 ]}\nobreak
\hyperdef{L}{X85FB9A8E7D9D5BD9}{}
{\noindent\textcolor{FuncColor}{$\triangleright$\enspace\texttt{FGAbElement({\mdseries\slshape A, v})\index{FGAbElement@\texttt{FGAbElement}!for IsFGAb, IsVector}
\label{FGAbElement:for IsFGAb, IsVector}
}\hfill{\scriptsize (operation)}}\\
\textbf{\indent Returns:\ }
a \texttt{FGAbElement} in \texttt{A} 



 \texttt{A} is a \texttt{FGAb} and \texttt{v} is a vector. }

 

\subsection{\textcolor{Chapter }{IsFGAbElement (for IsAdditiveElementWithInverse and IsAdditiveElementWithZero and IsExtLElement and IsExtRElement)}}
\logpage{[ 2, 1, 2 ]}\nobreak
\label{`IsFGAbElement`}
\hyperdef{L}{X78C70AB87C8D345C}{}
{\noindent\textcolor{FuncColor}{$\triangleright$\enspace\texttt{IsFGAbElement({\mdseries\slshape v})\index{IsFGAbElement@\texttt{IsFGAbElement}!for IsAdditiveElementWithInverse and IsAdditiveElementWithZero and IsExtLElement and IsExtRElement}
\label{IsFGAbElement:for IsAdditiveElementWithInverse and IsAdditiveElementWithZero and IsExtLElement and IsExtRElement}
}\hfill{\scriptsize (Category)}}\\
\noindent\textcolor{FuncColor}{$\triangleright$\enspace\texttt{IsFGAbElementCollection({\mdseries\slshape x})\index{IsFGAbElementCollection@\texttt{IsFGAbElementCollection}}
\label{IsFGAbElementCollection}
}\hfill{\scriptsize (Collection)}}\\
\noindent\textcolor{FuncColor}{$\triangleright$\enspace\texttt{IsFGAbElementCollColl({\mdseries\slshape x})\index{IsFGAbElementCollColl@\texttt{IsFGAbElementCollColl}}
\label{IsFGAbElementCollColl}
}\hfill{\scriptsize (Collection)}}\\
\textbf{\indent Returns:\ }
\texttt{true} or \texttt{false} 



 

 }

 }

 
\section{\textcolor{Chapter }{Additive Group}}\label{Chapter_`FGAb`_Section_Additive_Group}
\logpage{[ 2, 2, 0 ]}
\hyperdef{L}{X85E244B5797D4BBD}{}
{
  

 

\subsection{\textcolor{Chapter }{IsFGAb (for IsAdditiveGroup and IsFGAbElementCollection)}}
\logpage{[ 2, 2, 1 ]}\nobreak
\hyperdef{L}{X812F2E4D7EE2725A}{}
{\noindent\textcolor{FuncColor}{$\triangleright$\enspace\texttt{IsFGAb({\mdseries\slshape A})\index{IsFGAb@\texttt{IsFGAb}!for IsAdditiveGroup and IsFGAbElementCollection}
\label{IsFGAb:for IsAdditiveGroup and IsFGAbElementCollection}
}\hfill{\scriptsize (Category)}}\\
\textbf{\indent Returns:\ }
\texttt{true} or \texttt{false} 



 \texttt{IsFGAb} is the supercategory of \texttt{IsAdditiveGroup and IsFGAbElementCollection}. }

 

\subsection{\textcolor{Chapter }{UnderlyingMatrix (for IsFGAb and IsFGAbRep)}}
\logpage{[ 2, 2, 2 ]}\nobreak
\hyperdef{L}{X82A67F3784ED5DDC}{}
{\noindent\textcolor{FuncColor}{$\triangleright$\enspace\texttt{UnderlyingMatrix({\mdseries\slshape A})\index{UnderlyingMatrix@\texttt{UnderlyingMatrix}!for IsFGAb and IsFGAbRep}
\label{UnderlyingMatrix:for IsFGAb and IsFGAbRep}
}\hfill{\scriptsize (attribute)}}\\
\textbf{\indent Returns:\ }
the underlying matirx of \texttt{A} 



 

 }

 

\subsection{\textcolor{Chapter }{SimplifiedFGAb (for IsFGAb and IsFGAbRep)}}
\logpage{[ 2, 2, 3 ]}\nobreak
\label{`SimplifiedFGAb`}
\hyperdef{L}{X78580BDE798175FE}{}
{\noindent\textcolor{FuncColor}{$\triangleright$\enspace\texttt{SimplifiedFGAb({\mdseries\slshape A})\index{SimplifiedFGAb@\texttt{SimplifiedFGAb}!for IsFGAb and IsFGAbRep}
\label{SimplifiedFGAb:for IsFGAb and IsFGAbRep}
}\hfill{\scriptsize (attribute)}}\\
\noindent\textcolor{FuncColor}{$\triangleright$\enspace\texttt{ToSimplifiedFGAb({\mdseries\slshape A})\index{ToSimplifiedFGAb@\texttt{ToSimplifiedFGAb}!for IsFGAb and IsFGAbRep}
\label{ToSimplifiedFGAb:for IsFGAb and IsFGAbRep}
}\hfill{\scriptsize (attribute)}}\\
\noindent\textcolor{FuncColor}{$\triangleright$\enspace\texttt{FromSimplifiedFGAb({\mdseries\slshape A})\index{FromSimplifiedFGAb@\texttt{FromSimplifiedFGAb}!for IsFGAb and IsFGAbRep}
\label{FromSimplifiedFGAb:for IsFGAb and IsFGAbRep}
}\hfill{\scriptsize (attribute)}}\\
\textbf{\indent Returns:\ }
a \texttt{FGAb} 



 We can make a Schmit decompostion \texttt{UnderlyingMatrix(A)} to get a diagonal matrix, and discard redundant dimensions. The input must be
a output of \texttt{SimplifiedFGAb}. The input must be a output of \texttt{SimplifiedFGAb}. }

 

\subsection{\textcolor{Chapter }{FGAb (for IsVector)}}
\logpage{[ 2, 2, 4 ]}\nobreak
\hyperdef{L}{X8218CF358690CA7A}{}
{\noindent\textcolor{FuncColor}{$\triangleright$\enspace\texttt{FGAb({\mdseries\slshape m})\index{FGAb@\texttt{FGAb}!for IsVector}
\label{FGAb:for IsVector}
}\hfill{\scriptsize (operation)}}\\
\textbf{\indent Returns:\ }
the \texttt{FGAb} 



 \texttt{m} can be either a \texttt{MatrixObj} or a naive matirx. For example, if you want a $\mathbb{Z}_2\oplus \mathbb{Z}_4$, use \texttt{FGAb([[2, 0], [0, 4]])}. }

 

\subsection{\textcolor{Chapter }{FreeFGAb (for IsInt)}}
\logpage{[ 2, 2, 5 ]}\nobreak
\hyperdef{L}{X83E9554D87ADDB5C}{}
{\noindent\textcolor{FuncColor}{$\triangleright$\enspace\texttt{FreeFGAb({\mdseries\slshape n})\index{FreeFGAb@\texttt{FreeFGAb}!for IsInt}
\label{FreeFGAb:for IsInt}
}\hfill{\scriptsize (operation)}}\\
\textbf{\indent Returns:\ }
the free \texttt{FGAb}, i.e. $\mathbb{Z}^n$ 



 

 }

 

\subsection{\textcolor{Chapter }{SubFGAb (for IsFGAb, IsList)}}
\logpage{[ 2, 2, 6 ]}\nobreak
\label{`SubFGAb`}
\hyperdef{L}{X8504AC788030B079}{}
{\noindent\textcolor{FuncColor}{$\triangleright$\enspace\texttt{SubFGAb({\mdseries\slshape A, gens})\index{SubFGAb@\texttt{SubFGAb}!for IsFGAb, IsList}
\label{SubFGAb:for IsFGAb, IsList}
}\hfill{\scriptsize (operation)}}\\
\noindent\textcolor{FuncColor}{$\triangleright$\enspace\texttt{FromSubFGAb({\mdseries\slshape A})\index{FromSubFGAb@\texttt{FromSubFGAb}!for IsFGAb and IsFGAbRep}
\label{FromSubFGAb:for IsFGAb and IsFGAbRep}
}\hfill{\scriptsize (attribute)}}\\
\textbf{\indent Returns:\ }
the \texttt{SubFGAb} generated by \texttt{gens} 



 

 The input \texttt{A} should be the output of \texttt{SubFGAb}. }

 

\subsection{\textcolor{Chapter }{IndependentGeneratorsOfAdditiveGroup (for IsAdditiveGroup)}}
\logpage{[ 2, 2, 7 ]}\nobreak
\hyperdef{L}{X81A7960B80AEBBDF}{}
{\noindent\textcolor{FuncColor}{$\triangleright$\enspace\texttt{IndependentGeneratorsOfAdditiveGroup({\mdseries\slshape A})\index{IndependentGeneratorsOfAdditiveGroup@\texttt{Independent}\-\texttt{Generators}\-\texttt{Of}\-\texttt{Additive}\-\texttt{Group}!for IsAdditiveGroup}
\label{IndependentGeneratorsOfAdditiveGroup:for IsAdditiveGroup}
}\hfill{\scriptsize (attribute)}}\\
\textbf{\indent Returns:\ }
a list of generators 



 As \texttt{SimplifiedFGAb(a)} gives the \texttt{}smallest'' representation of \texttt{A}, we can use it as a \texttt{IndependentGeneratorsOfAdditiveGroup}. }

 

\subsection{\textcolor{Chapter }{FGAbDirectSum (for IsFGAb, IsFGAb)}}
\logpage{[ 2, 2, 8 ]}\nobreak
\label{`FGAbDirectSum`}
\hyperdef{L}{X7C9D77E27FA5601A}{}
{\noindent\textcolor{FuncColor}{$\triangleright$\enspace\texttt{FGAbDirectSum({\mdseries\slshape A1, A2})\index{FGAbDirectSum@\texttt{FGAbDirectSum}!for IsFGAb, IsFGAb}
\label{FGAbDirectSum:for IsFGAb, IsFGAb}
}\hfill{\scriptsize (operation)}}\\
\noindent\textcolor{FuncColor}{$\triangleright$\enspace\texttt{FGAbDirectSum({\mdseries\slshape listA})\index{FGAbDirectSum@\texttt{FGAbDirectSum}!for IsList}
\label{FGAbDirectSum:for IsList}
}\hfill{\scriptsize (operation)}}\\
\noindent\textcolor{FuncColor}{$\triangleright$\enspace\texttt{FGAbDirectSumInfo({\mdseries\slshape A})\index{FGAbDirectSumInfo@\texttt{FGAbDirectSumInfo}!for IsFGAb}
\label{FGAbDirectSumInfo:for IsFGAb}
}\hfill{\scriptsize (attribute)}}\\
\noindent\textcolor{FuncColor}{$\triangleright$\enspace\texttt{Embedding({\mdseries\slshape A, n})\index{Embedding@\texttt{Embedding}!for IsFGAb and IsFGAbRep and HasFGAbDirectSumInfo, IsPosInt}
\label{Embedding:for IsFGAb and IsFGAbRep and HasFGAbDirectSumInfo, IsPosInt}
}\hfill{\scriptsize (method)}}\\
\noindent\textcolor{FuncColor}{$\triangleright$\enspace\texttt{Projection({\mdseries\slshape A, n})\index{Projection@\texttt{Projection}!for IsFGAb and IsFGAbRep and HasFGAbDirectSumInfo, IsPosInt}
\label{Projection:for IsFGAb and IsFGAbRep and HasFGAbDirectSumInfo, IsPosInt}
}\hfill{\scriptsize (method)}}\\


 

 \texttt{FGAbDirectSumInfo} stored a record with \texttt{FGAb}, \texttt{embeddings}, \texttt{projections}. \texttt{FGAb} is a list of \texttt{FGAb}, \texttt{embeddings} and \texttt{projections} are lists of canonical maps. \texttt{Embedding} can be used. \texttt{Projection} can be used. }

 

\subsection{\textcolor{Chapter }{GeneratorsOfAdditiveGroup (for IsFGAb and IsFGAbRep)}}
\logpage{[ 2, 2, 9 ]}\nobreak
\hyperdef{L}{X86EC1DD67A24122F}{}
{\noindent\textcolor{FuncColor}{$\triangleright$\enspace\texttt{GeneratorsOfAdditiveGroup({\mdseries\slshape A})\index{GeneratorsOfAdditiveGroup@\texttt{GeneratorsOfAdditiveGroup}!for IsFGAb and IsFGAbRep}
\label{GeneratorsOfAdditiveGroup:for IsFGAb and IsFGAbRep}
}\hfill{\scriptsize (method)}}\\
\textbf{\indent Returns:\ }
a list of generators 



 

 }

 

\subsection{\textcolor{Chapter }{Intersection2 (for IsAdditiveGroup and IsFGAbElementCollection, IsAdditiveGroup and IsFGAbElementCollection)}}
\logpage{[ 2, 2, 10 ]}\nobreak
\hyperdef{L}{X7A40CD9D7BDD45BB}{}
{\noindent\textcolor{FuncColor}{$\triangleright$\enspace\texttt{Intersection2({\mdseries\slshape A1, A2})\index{Intersection2@\texttt{Intersection2}!for IsAdditiveGroup and IsFGAbElementCollection, IsAdditiveGroup and IsFGAbElementCollection}
\label{Intersection2:for IsAdditiveGroup and IsFGAbElementCollection, IsAdditiveGroup and IsFGAbElementCollection}
}\hfill{\scriptsize (method)}}\\
\textbf{\indent Returns:\ }
a \texttt{SubFGAb} 



 You can use just \texttt{Intersection}. \texttt{A1} and \texttt{A2} should both \texttt{HasParent}, or \texttt{Parent(A1)} is \texttt{A2}, or \texttt{Parent(A2)} is \texttt{A1}. }

 

\subsection{\textcolor{Chapter }{\texttt{\symbol{92}}+ (for IsAdditiveGroup and IsFGAbElementCollection, IsAdditiveGroup and IsFGAbElementCollection)}}
\logpage{[ 2, 2, 11 ]}\nobreak
\hyperdef{L}{X7809D15B84AECCB2}{}
{\noindent\textcolor{FuncColor}{$\triangleright$\enspace\texttt{\texttt{\symbol{92}}+({\mdseries\slshape A1, A2})\index{\texttt{\symbol{92}}+@\texttt{\texttt{\symbol{92}}+}!for IsAdditiveGroup and IsFGAbElementCollection, IsAdditiveGroup and IsFGAbElementCollection}
\label{bSlash+:for IsAdditiveGroup and IsFGAbElementCollection, IsAdditiveGroup and IsFGAbElementCollection}
}\hfill{\scriptsize (method)}}\\
\textbf{\indent Returns:\ }
a \texttt{SubFGAb} \texttt{A1 + A2} 



 \texttt{A1} and \texttt{A2} should both \texttt{HasParent}, or \texttt{Parent(A1)} is \texttt{A2}, or \texttt{Parent(A2)} is \texttt{A1}. }

 

\subsection{\textcolor{Chapter }{IsTrivial (for IsFGAb and IsFGAbRep)}}
\logpage{[ 2, 2, 12 ]}\nobreak
\hyperdef{L}{X7E587C87785E33FD}{}
{\noindent\textcolor{FuncColor}{$\triangleright$\enspace\texttt{IsTrivial({\mdseries\slshape A})\index{IsTrivial@\texttt{IsTrivial}!for IsFGAb and IsFGAbRep}
\label{IsTrivial:for IsFGAb and IsFGAbRep}
}\hfill{\scriptsize (method)}}\\


 Determine whether A is a trivial additive group. }

 

\subsection{\textcolor{Chapter }{\texttt{\symbol{92}}= (for IsAdditiveGroup and IsFGAbElementCollection, IsAdditiveGroup and IsFGAbElementCollection)}}
\logpage{[ 2, 2, 13 ]}\nobreak
\hyperdef{L}{X7A3FEFDF7BE387E7}{}
{\noindent\textcolor{FuncColor}{$\triangleright$\enspace\texttt{\texttt{\symbol{92}}=({\mdseries\slshape A1, A2})\index{\texttt{\symbol{92}}=@\texttt{\texttt{\symbol{92}}=}!for IsAdditiveGroup and IsFGAbElementCollection, IsAdditiveGroup and IsFGAbElementCollection}
\label{bSlash=:for IsAdditiveGroup and IsFGAbElementCollection, IsAdditiveGroup and IsFGAbElementCollection}
}\hfill{\scriptsize (method)}}\\


 Determine whether \texttt{A1 = A2}. \texttt{A1} and \texttt{A2} should both \texttt{HasParent}, or \texttt{Parent(A1)} is \texttt{A2}, or \texttt{Parent(A2)} is \texttt{A1}. }

 }

 }

   
\chapter{\textcolor{Chapter }{\texttt{FGAbHomomorphism}}}\label{Chapter_`FGAbHomomorphism`}
\logpage{[ 3, 0, 0 ]}
\hyperdef{L}{X7E7B90F37DDBF289}{}
{
  

 
\section{\textcolor{Chapter }{\texttt{FGAbHomomorphism}}}\label{Chapter_`FGAbHomomorphism`_Section_`FGAbHomomorphism`}
\logpage{[ 3, 1, 0 ]}
\hyperdef{L}{X7E7B90F37DDBF289}{}
{
  

 
\subsection{\textcolor{Chapter }{\texttt{Image}, \texttt{Images}}}\label{Chapter_`FGAbHomomorphism`_Section_`FGAbHomomorphism`_Subsection_`Image`,_`Images`}
\logpage{[ 3, 1, 1 ]}
\hyperdef{L}{X87F942ED85370A8A}{}
{
  

 \texttt{Image} and \texttt{Images} can be used directly even for a \texttt{SubFGAb}. 

 }

 
\subsection{\textcolor{Chapter }{\texttt{PreImage}, \texttt{PreImages}}}\label{Chapter_`FGAbHomomorphism`_Section_`FGAbHomomorphism`_Subsection_`PreImage`,_`PreImages`}
\logpage{[ 3, 1, 2 ]}
\hyperdef{L}{X7B8B0F7A78067FF4}{}
{
  

 \texttt{PreImage} and \texttt{PreImages} only works for a \texttt{SubFGAb}, since the preimage of an element is not a \texttt{FGAb} but a additive coset, our method can't apply. However, \texttt{PreImagesRepresentative} is available. 

 }

 

\subsection{\textcolor{Chapter }{IsFGAbHomomorphism (for IsAdditiveGroupGeneralMapping)}}
\logpage{[ 3, 1, 3 ]}\nobreak
\hyperdef{L}{X80E6E0F67EB4EE86}{}
{\noindent\textcolor{FuncColor}{$\triangleright$\enspace\texttt{IsFGAbHomomorphism({\mdseries\slshape m})\index{IsFGAbHomomorphism@\texttt{IsFGAbHomomorphism}!for IsAdditiveGroupGeneralMapping}
\label{IsFGAbHomomorphism:for IsAdditiveGroupGeneralMapping}
}\hfill{\scriptsize (Category)}}\\
\textbf{\indent Returns:\ }
\texttt{true} or \texttt{false} 



 

 }

 

\subsection{\textcolor{Chapter }{FGAbHomomorphism (for IsFGAb, IsFGAb, IsMatrixOrMatrixObj)}}
\logpage{[ 3, 1, 4 ]}\nobreak
\label{`FGAbHomomorphism`}
\hyperdef{L}{X7E77D0B08589328A}{}
{\noindent\textcolor{FuncColor}{$\triangleright$\enspace\texttt{FGAbHomomorphism({\mdseries\slshape A, B, m})\index{FGAbHomomorphism@\texttt{FGAbHomomorphism}!for IsFGAb, IsFGAb, IsMatrixOrMatrixObj}
\label{FGAbHomomorphism:for IsFGAb, IsFGAb, IsMatrixOrMatrixObj}
}\hfill{\scriptsize (operation)}}\\
\noindent\textcolor{FuncColor}{$\triangleright$\enspace\texttt{FGAbHomomorphismNC({\mdseries\slshape A, B, m})\index{FGAbHomomorphismNC@\texttt{FGAbHomomorphismNC}!for IsFGAb, IsFGAb, IsMatrixOrMatrixObj}
\label{FGAbHomomorphismNC:for IsFGAb, IsFGAb, IsMatrixOrMatrixObj}
}\hfill{\scriptsize (operation)}}\\
\noindent\textcolor{FuncColor}{$\triangleright$\enspace\texttt{FGAbHomomorphism({\mdseries\slshape A, B, m, n})\index{FGAbHomomorphism@\texttt{FGAbHomomorphism}!for IsFGAb, IsFGAb, IsMatrixOrMatrixObj, IsMatrixOrMatrixObj}
\label{FGAbHomomorphism:for IsFGAb, IsFGAb, IsMatrixOrMatrixObj, IsMatrixOrMatrixObj}
}\hfill{\scriptsize (operation)}}\\
\noindent\textcolor{FuncColor}{$\triangleright$\enspace\texttt{FGAbHomomorphismNC({\mdseries\slshape A, B, m, n})\index{FGAbHomomorphismNC@\texttt{FGAbHomomorphismNC}!for IsFGAb, IsFGAb, IsMatrixOrMatrixObj, IsMatrixOrMatrixObj}
\label{FGAbHomomorphismNC:for IsFGAb, IsFGAb, IsMatrixOrMatrixObj, IsMatrixOrMatrixObj}
}\hfill{\scriptsize (operation)}}\\
\textbf{\indent Returns:\ }
a \texttt{FGAbHomomorphism} 



 \texttt{A} and \texttt{B} are \texttt{FGAb} as domain and codomain. Let \texttt{A} be $A_1\to A_0$, and \texttt{B} be $B_1\to B_0$, \texttt{m} is a matrix representing a map $A_0\to B_0$. \texttt{n} is a matrix representing a map $A_1\to B_1$. \texttt{n} is actually determined by \texttt{m} up to a chain homotopy. Therefore, we don't need to give \texttt{n} explicitly. 

 }

 

\subsection{\textcolor{Chapter }{UnderlyingMatrix (for IsFGAbHomomorphism and IsFGAbHomomorphismRep)}}
\logpage{[ 3, 1, 5 ]}\nobreak
\hyperdef{L}{X7A97BF7B7F3E9A3B}{}
{\noindent\textcolor{FuncColor}{$\triangleright$\enspace\texttt{UnderlyingMatrix({\mdseries\slshape A})\index{UnderlyingMatrix@\texttt{UnderlyingMatrix}!for IsFGAbHomomorphism and IsFGAbHomomorphismRep}
\label{UnderlyingMatrix:for IsFGAbHomomorphism and IsFGAbHomomorphismRep}
}\hfill{\scriptsize (attribute)}}\\
\textbf{\indent Returns:\ }
the underlying matirx $A_0\to B_0$ of a \texttt{FGAbHomomorphism} 



 

 }

 

\subsection{\textcolor{Chapter }{LiftingMatrix (for IsFGAbHomomorphism and IsFGAbHomomorphismRep)}}
\logpage{[ 3, 1, 6 ]}\nobreak
\hyperdef{L}{X7E90BCEA7AD9EDEB}{}
{\noindent\textcolor{FuncColor}{$\triangleright$\enspace\texttt{LiftingMatrix({\mdseries\slshape A})\index{LiftingMatrix@\texttt{LiftingMatrix}!for IsFGAbHomomorphism and IsFGAbHomomorphismRep}
\label{LiftingMatrix:for IsFGAbHomomorphism and IsFGAbHomomorphismRep}
}\hfill{\scriptsize (attribute)}}\\
\textbf{\indent Returns:\ }
the underlying matirx $A_1\to B_1$ of a \texttt{FGAbHomomorphism} 



 

 }

 }

 
\section{\textcolor{Chapter }{limits and colimits}}\label{Chapter_`FGAbHomomorphism`_Section_limits_and_colimits}
\logpage{[ 3, 2, 0 ]}
\hyperdef{L}{X857CF4D487158F26}{}
{
  

\subsection{\textcolor{Chapter }{FGAbKernel (for IsFGAbHomomorphism and IsFGAbHomomorphismRep)}}
\logpage{[ 3, 2, 1 ]}\nobreak
\label{`FGAbKernel`}
\hyperdef{L}{X7A38EC84839A5408}{}
{\noindent\textcolor{FuncColor}{$\triangleright$\enspace\texttt{FGAbKernel({\mdseries\slshape f})\index{FGAbKernel@\texttt{FGAbKernel}!for IsFGAbHomomorphism and IsFGAbHomomorphismRep}
\label{FGAbKernel:for IsFGAbHomomorphism and IsFGAbHomomorphismRep}
}\hfill{\scriptsize (attribute)}}\\
\noindent\textcolor{FuncColor}{$\triangleright$\enspace\texttt{FromFGAbKernel({\mdseries\slshape A})\index{FromFGAbKernel@\texttt{FromFGAbKernel}!for IsFGAb and IsFGAbRep}
\label{FromFGAbKernel:for IsFGAb and IsFGAbRep}
}\hfill{\scriptsize (attribute)}}\\
\noindent\textcolor{FuncColor}{$\triangleright$\enspace\texttt{UnderlyingFGAbKernel({\mdseries\slshape A})\index{UnderlyingFGAbKernel@\texttt{UnderlyingFGAbKernel}!for IsAdditiveGroup}
\label{UnderlyingFGAbKernel:for IsAdditiveGroup}
}\hfill{\scriptsize (attribute)}}\\
\noindent\textcolor{FuncColor}{$\triangleright$\enspace\texttt{KernelOfAdditiveGeneralMapping({\mdseries\slshape f})\index{KernelOfAdditiveGeneralMapping@\texttt{KernelOfAdditiveGeneralMapping}!for IsFGAbHomomorphism and IsFGAbHomomorphismRep}
\label{KernelOfAdditiveGeneralMapping:for IsFGAbHomomorphism and IsFGAbHomomorphismRep}
}\hfill{\scriptsize (method)}}\\


 \texttt{FGAbKernel(f)} is a \texttt{FGAb} equipped with \texttt{FromFGAbKernel}. The input of \texttt{FromFGAbKernel} must be a output of \texttt{FGAbKernel}. 

 You can use just \texttt{Kernel(f)} to get the same result as \texttt{KernelOfAdditiveGeneralMapping(f)}. \texttt{Kernel(f)} has an attribute \texttt{UnderlyingFGAbKernel(A)} which is a output of \texttt{FGAbKernel}. 

 \texttt{Kernel(f)} is what you should use. Because \texttt{Kernel(f)} is actually a subset of \texttt{Range(f)}, in particular, a SubFGAb. However, \texttt{FGAbKernel(f)} is just a \texttt{FGAb} equipped with a canonical map. \texttt{FGAbKernel} should be treated as an inner\texttt{\symbol{45}}built function, which is used
to get \texttt{kernel(f)}. }

 

\subsection{\textcolor{Chapter }{FGAbCokernel (for IsFGAbHomomorphism and IsFGAbHomomorphismRep)}}
\logpage{[ 3, 2, 2 ]}\nobreak
\label{`FGAbCokernel`}
\hyperdef{L}{X7E7E8F57805CB1DF}{}
{\noindent\textcolor{FuncColor}{$\triangleright$\enspace\texttt{FGAbCokernel({\mdseries\slshape f})\index{FGAbCokernel@\texttt{FGAbCokernel}!for IsFGAbHomomorphism and IsFGAbHomomorphismRep}
\label{FGAbCokernel:for IsFGAbHomomorphism and IsFGAbHomomorphismRep}
}\hfill{\scriptsize (attribute)}}\\
\noindent\textcolor{FuncColor}{$\triangleright$\enspace\texttt{ToFGAbCokernel({\mdseries\slshape A})\index{ToFGAbCokernel@\texttt{ToFGAbCokernel}!for IsFGAb and IsFGAbRep}
\label{ToFGAbCokernel:for IsFGAb and IsFGAbRep}
}\hfill{\scriptsize (attribute)}}\\


 \texttt{FGAbKernel(f)} is a \texttt{FGAb} equipped with \texttt{ToFGAbCokernel}. The input of \texttt{FromFGAbCokernel} must be a output of \texttt{FGAbCokernel}. }

 

\subsection{\textcolor{Chapter }{FGAbPullback (for IsFGAbHomomorphism and IsFGAbHomomorphismRep, IsFGAbHomomorphism and IsFGAbHomomorphismRep)}}
\logpage{[ 3, 2, 3 ]}\nobreak
\label{`FGAbPullback`}
\hyperdef{L}{X87A48093864D9E0E}{}
{\noindent\textcolor{FuncColor}{$\triangleright$\enspace\texttt{FGAbPullback({\mdseries\slshape f1, f2})\index{FGAbPullback@\texttt{FGAbPullback}!for IsFGAbHomomorphism and IsFGAbHomomorphismRep, IsFGAbHomomorphism and IsFGAbHomomorphismRep}
\label{FGAbPullback:for IsFGAbHomomorphism and IsFGAbHomomorphismRep, IsFGAbHomomorphism and IsFGAbHomomorphismRep}
}\hfill{\scriptsize (operation)}}\\
\noindent\textcolor{FuncColor}{$\triangleright$\enspace\texttt{FromFGAbPullback({\mdseries\slshape A, n})\index{FromFGAbPullback@\texttt{FromFGAbPullback}!for IsFGAb and IsFGAbRep, IsInt}
\label{FromFGAbPullback:for IsFGAb and IsFGAbRep, IsInt}
}\hfill{\scriptsize (operation)}}\\


 \texttt{FGAbPullback(f1, f2)} is a \texttt{FGAb} equipped with \texttt{FromFGAbPullback}. \texttt{FromFGAbPullback(A, n)} is the projection to \texttt{Source(f1)} or \texttt{Source(f2)}. }

 

\subsection{\textcolor{Chapter }{FGAbPushout (for IsFGAbHomomorphism and IsFGAbHomomorphismRep, IsFGAbHomomorphism and IsFGAbHomomorphismRep)}}
\logpage{[ 3, 2, 4 ]}\nobreak
\label{`FGAbPushout`}
\hyperdef{L}{X7C87F9E685850A2D}{}
{\noindent\textcolor{FuncColor}{$\triangleright$\enspace\texttt{FGAbPushout({\mdseries\slshape f1, f2})\index{FGAbPushout@\texttt{FGAbPushout}!for IsFGAbHomomorphism and IsFGAbHomomorphismRep, IsFGAbHomomorphism and IsFGAbHomomorphismRep}
\label{FGAbPushout:for IsFGAbHomomorphism and IsFGAbHomomorphismRep, IsFGAbHomomorphism and IsFGAbHomomorphismRep}
}\hfill{\scriptsize (operation)}}\\
\noindent\textcolor{FuncColor}{$\triangleright$\enspace\texttt{ToFGAbPushout({\mdseries\slshape A, n})\index{ToFGAbPushout@\texttt{ToFGAbPushout}!for IsFGAb and IsFGAbRep, IsInt}
\label{ToFGAbPushout:for IsFGAb and IsFGAbRep, IsInt}
}\hfill{\scriptsize (operation)}}\\


 \texttt{FGAbPushout(f1, f2)} is a \texttt{FGAb} equipped with \texttt{ToFGAbPushout}. \texttt{ToFGAbPushout(A, n)} is the projection to \texttt{Range(f1)} or \texttt{Range(f2)}. }

 }

 }

 \def\indexname{Index\logpage{[ "Ind", 0, 0 ]}
\hyperdef{L}{X83A0356F839C696F}{}
}

\cleardoublepage
\phantomsection
\addcontentsline{toc}{chapter}{Index}


\printindex

\immediate\write\pagenrlog{["Ind", 0, 0], \arabic{page},}
\immediate\write\pagenrlog{["Ind", 0, 0], \arabic{page},}
\immediate\write\pagenrlog{["Ind", 0, 0], \arabic{page},}
\immediate\write\pagenrlog{["Ind", 0, 0], \arabic{page},}
\immediate\write\pagenrlog{["Ind", 0, 0], \arabic{page},}
\immediate\write\pagenrlog{["Ind", 0, 0], \arabic{page},}
\immediate\write\pagenrlog{["Ind", 0, 0], \arabic{page},}
\immediate\write\pagenrlog{["Ind", 0, 0], \arabic{page},}
\immediate\write\pagenrlog{["Ind", 0, 0], \arabic{page},}
\immediate\write\pagenrlog{["Ind", 0, 0], \arabic{page},}
\immediate\write\pagenrlog{["Ind", 0, 0], \arabic{page},}
\immediate\write\pagenrlog{["Ind", 0, 0], \arabic{page},}
\immediate\write\pagenrlog{["Ind", 0, 0], \arabic{page},}
\immediate\write\pagenrlog{["Ind", 0, 0], \arabic{page},}
\immediate\write\pagenrlog{["Ind", 0, 0], \arabic{page},}
\immediate\write\pagenrlog{["Ind", 0, 0], \arabic{page},}
\immediate\write\pagenrlog{["Ind", 0, 0], \arabic{page},}
\immediate\write\pagenrlog{["Ind", 0, 0], \arabic{page},}
\immediate\write\pagenrlog{["Ind", 0, 0], \arabic{page},}
\immediate\write\pagenrlog{["Ind", 0, 0], \arabic{page},}
\immediate\write\pagenrlog{["Ind", 0, 0], \arabic{page},}
\immediate\write\pagenrlog{["Ind", 0, 0], \arabic{page},}
\immediate\write\pagenrlog{["Ind", 0, 0], \arabic{page},}
\immediate\write\pagenrlog{["Ind", 0, 0], \arabic{page},}
\immediate\write\pagenrlog{["Ind", 0, 0], \arabic{page},}
\immediate\write\pagenrlog{["Ind", 0, 0], \arabic{page},}
\immediate\write\pagenrlog{["Ind", 0, 0], \arabic{page},}
\immediate\write\pagenrlog{["Ind", 0, 0], \arabic{page},}
\immediate\write\pagenrlog{["Ind", 0, 0], \arabic{page},}
\immediate\write\pagenrlog{["Ind", 0, 0], \arabic{page},}
\immediate\write\pagenrlog{["Ind", 0, 0], \arabic{page},}
\immediate\write\pagenrlog{["Ind", 0, 0], \arabic{page},}
\immediate\write\pagenrlog{["Ind", 0, 0], \arabic{page},}
\immediate\write\pagenrlog{["Ind", 0, 0], \arabic{page},}
\immediate\write\pagenrlog{["Ind", 0, 0], \arabic{page},}
\immediate\write\pagenrlog{["Ind", 0, 0], \arabic{page},}
\newpage
\immediate\write\pagenrlog{["End"], \arabic{page}];}
\immediate\closeout\pagenrlog
\end{document}
